\rsection{Soziale Kognition}
\slide{Cheater-Detection (Betrüger-Erkennung)}{
    \begin{itemize}
        \item Erklärung: Menschen besitzen spezialisierte kognitive Mechanismen (oder nutzen Alltagswissen), um Verstöße gegen soziale Regeln oder Verträge zu erkennen (Cheater-Detection).
        \item Universalität: Dieser Effekt ist unabhängig von der Kultur und der genauen Formulierung.
    \end{itemize}
}
\slide{Vorurteile, Stereotype \& Einstellungen}{
    \begin{itemize}
        \item \b{Stereotyp:} Verallgemeinernde Zuschreibung von Eigenschaften auf alle Mitglieder einer Gruppe, ohne individuelle Unterschiede zu beachten.
        \item \b{Vorurteil:} Eine negative Bewertung und ablehnende Haltung gegenüber Personen allein aufgrund ihrer Gruppenzugehörigkeit.
        \item \b{Diskriminierung:} Das tatsächliche Verhalten, das Personen aufgrund von Stereotypen oder Vorurteilen benachteiligt.
        \item \b{Funktion:} Dient der kognitiven Entlastung durch Komplexitätsreduktion (Kategorisierung).
        \item \b{Einstellung vs. Verhalten:} Es gibt oft eine Diskrepanz zwischen der berichteten Einstellung (z. B. Fairness in Umfragen) und dem tatsächlichen Verhalten (z. B. egoistische Aufteilung von Aufgaben).
    \end{itemize}
}
\slide{Theory of Mind (ToM)}{
    \begin{itemize}
        \item Definition: Das Wissen darüber, dass andere Personen eigene Motivationen, Gefühle und Wissensstände haben, die von den eigenen oder der Realität abweichen können.
        \item \b{Kernaspekt:} Verständnis, dass der Zustand der Welt nicht zwingend mit dem Wissensstand einer anderen Person übereinstimmen muss (z. B. "Sally-Anne-Test").
        \item Komponenten:
        \begin{itemize}
            \item Arbeitsgedächtnis: Gleichzeitiges Vorhalten mehrerer Zustände (Realität vs. Glaube des anderen).
            \item Exekutivfunktionen: Fähigkeit zum Perspektivwechsel.
            \item Sprache: Nutzung psychologischer Begriffe (z. B. "glauben", "wissen").
        \end{itemize}
    \end{itemize}
}
\slide{Kognition in Gruppen}{
    \begin{itemize}
        \item \b{Alignment:} Gesprächspartner gleichen sich an und definieren gemeinsame Sprachmuster, die spezifisch für ihre Interaktion sind.
        \item \b{Gruppen-Dynamik:}
        \begin{itemize}
            \item Kleine Gruppen: Kommunikative Sprecher haben den größten Einfluss.
            \item Große Gruppen: Dominante Sprecher haben den größten Einfluss.
        \end{itemize}
        \item \b{Effektivität:} Gruppen arbeiten effektiver, wenn sie externe Hilfsmittel (Notizen, Artefakte) nutzen.
        \item Fazit: Kognition ist untrennbar mit dem sozialen Kontext verbunden; die Grenzen zur Sozialpsychologie sind fließend.
    \end{itemize}
}